\documentclass[conference]{IEEEtran}

\usepackage[T1]{fontenc}
\usepackage[utf8]{inputenc}
\usepackage{amsmath,amssymb,amsfonts}
\usepackage{graphicx}
\usepackage{booktabs}
\usepackage{listings}
\usepackage{xcolor}
\usepackage[hidelinks]{hyperref}

% Lean 4 listings setup (minimal; extend as needed)
\lstdefinelanguage{lean}{
  keywords={theorem,lemma,def,structure,instance,noncomputable,exists,fun,let,match,with,by,import,open,namespace,end},
  sensitive=true,
  comment=[l]{--},
  morecomment=[s]{/-}{-/},
}
\lstset{
  language=lean,
  % render the Unicode used in verbatim-quoted Lean statements under pdflatex
  literate={∃}{{$\exists$}}1 {ν}{{$\nu$}}1 {ℝ}{{$\mathbb{R}$}}1
           {₀}{{$_0$}}1 {𝔊}{{$\mathfrak{G}$}}1,
  basicstyle=\ttfamily\footnotesize,
  keywordstyle=\bfseries,
  commentstyle=\itshape\color{gray},
  breaklines=true,
  columns=fullflexible,
}

% Evidence cross-reference (draft-only; renders as a marginal tag, disabled for submission)
\newcommand{\evref}[1]{\textcolor{blue}{\footnotesize[#1]}}

\begin{document}

\title{Formalizing Leray--Hopf in Lean~4 with AI Agents:\\
Harness Design and Reliability Lessons}

\author{\IEEEauthorblockN{Anonymous Author(s)}}

\maketitle

\begin{abstract}
Formalization projects that run for weeks require safeguards beyond the
proof kernel. The kernel checks a theorem against its stated type, but it
does not choose the intended statement, preserve a public interface during
refactoring, or manage the agents and computing resources used to produce
the proof. We report a Lean~4 and mathlib formalization, with no project
axioms, of Leray--Hopf weak solution existence for the homogeneous
Navier--Stokes equations on $\mathbb{T}^3$ and $\mathbb{R}^3$. To our
knowledge, a July 2026 survey found no earlier existence theory for a
nonlinear fluid PDE that had been machine checked in Lean, Coq, or
Isabelle/HOL. During six weeks of development, a workflow defined by the
owner routed GitHub Issues through isolated Git worktrees and pull requests.
The harness routes
mathematical, Lean, interface, and resource checks according to the
artifact being changed. Four incidents show where
these checks detected defects and where they did not. The evidence method
links manuscript claims to redacted session records, repository history,
verification results, and measured costs.
\end{abstract}
% Abstract claims map: CLM-001 (result), CLM-002 (scale, lower bounds),
% CLM-006 (survey novelty), CLM-004/005 (incidents account),
% CLM-003 (billing reconciliation), CLM-007/009 (harness and orchestration).

\begin{IEEEkeywords}
formal verification, Lean 4, Navier--Stokes equations, Leray--Hopf weak solutions,
AI agents, large language models, proof engineering
\end{IEEEkeywords}

\section{Introduction}\label{sec:introduction}

Large language model agents can now sustain mathematical formalization
efforts spanning weeks and more than a hundred pull requests, rather
than single proof-search episodes, shifting the central reliability
question. A proof assistant's kernel guarantees that a \emph{proved}
declaration is correct given its stated type, but kernel checking is
silent on whether that type is the theorem the author actually
intended; whether a previously proved declaration silently disappeared
during a refactor; whether a build genuinely succeeded, versus being
reported successful over a stale disk state; and whether the project's
own development record is itself accurate. These statement-, process-,
and record-level failure modes are the load-bearing risk that remains
once proof-checking is presumed reliable.

We report on a formalization campaign that ran under exactly this
regime. Over roughly six weeks (2026-06-10 through the 2026-07-20
release day), the project consumed, on evidence collected so far, at
least 15.1M output tokens, an API-equivalent cost of \$6{,}972.83, and
199.1\,h of active engagement time under a five-minute inter-event gap
cap -- documented lower bounds, since coverage gaps in the session-log
record remain. % CLAIM: CLM-002
The campaign's central artifact is a kernel-only Lean~4~\cite{moura2021lean4}
formalization, built on mathlib~\cite{mathlib2020}, of existence of
Leray--Hopf weak solutions~\cite{leray1934,hopf1951} to the homogeneous,
incompressible Navier--Stokes equations, comprising two capstone existence
theorems -- one on the periodic 3-torus $\mathbb{T}^3$, one on
$\mathbb{R}^3$ -- whose \texttt{\#print axioms} output is exactly the
three standard Lean kernel axioms, with zero project-declared axioms and
no \texttt{sorryAx} on either theorem~\cite{lerayhopf2026repo}.
% CLAIM: CLM-001
To our knowledge, based on a literature and repository survey conducted in
July 2026 and summarized in the related-work discussion below, no
machine-checked existence theory for a nonlinear fluid PDE has previously
been reported in Lean, Coq, or Isabelle/HOL; we state this as a bound on
what our survey found, not as a proof of non-existence.\footnote{Search
terms, dates, and the nearest adjacent results are recorded in the
project's related-work notes and re-checked immediately before
submission.}
% CLAIM: CLM-006

Three literatures bound this work. Among AI-authored Lean projects,
TauCeti~\cite{tauceti2026}, incubated by the Lean FRO and the Mathlib
Initiative, pursues breadth over a thirteen-area roadmap under human
governance, but reports no accompanying paper, incident-level account,
or session-level evidence trail; we instead pursue depth against a
single frozen scope, with a named-incident case study as a central
artifact. Among formalized analysis adjacent to fluid PDE, mathlib's
Gagliardo--Nirenberg--Sobolev inequality~\cite{vandoorn2024gns} and a
De Giorgi--Nash--Moser elliptic regularity
formalization~\cite{armstrong2026degiorginashmoser} establish a
building-block inequality and a regularity theory respectively, neither
an existence theorem for a nonlinear PDE; the nearest precedent, an
AI-assisted formalization of well-posedness for the linear, kinetic
Vlasov equation~\cite{miller2026vlasov}, is single-author and
single-agent, with no multi-agent governance or incident structure.
Large-scale human-team formalizations such as the Odd Order
Theorem~\cite{gonthier2013oddorder} and the Liquid Tensor
Experiment~\cite{commelin2023abstraction} manage complexity through a
human team and a specification-driven blueprint; our setting substitutes
a single scientific owner governing agent-authored implementation, with
a frozen scope document as blueprint. Finally, the agentic and LLM
theorem-proving
literature~\cite{xin2025deepseekproverv2,yang2023leandojo,deepmind2025alphaproof,zhou2026leanatlas}
reports session- or benchmark-scale results, not a project-scale
process with recorded incidents; rising autonomous task
horizons~\cite{metr2025longtasks} and a study of agentic pull
requests~\cite{agenticprs2026security} motivate the reliability question
from outside formalization but do not instrument a single campaign to
this depth.

This paper makes three contributions:
\begin{enumerate}
\item A completed, kernel-only Lean~4 and mathlib formalization of
  Leray--Hopf weak-solution existence on $\mathbb{T}^3$ and $\mathbb{R}^3$,
  under an agent-orchestrated regime over roughly six weeks.
\item A workflow design -- statement-first development, independent
  adversarial review, and machine-checked verification gates -- with an
  evidence-backed account of what it caught and missed, drawn from named
  incidents rather than impressions.
\item An evidence methodology for AI-assisted research: session-log
  escrow with a manifest, machine-checked claim-to-evidence verification,
  and cost accounting reconciled against billing records.
\end{enumerate}

Section~\ref{sec:formalization} states the two capstone theorems.
Section~\ref{sec:workflow} describes the agent-orchestrated workflow.
Section~\ref{sec:incidents} presents four incidents where a build or
axiom check alone was insufficient to guarantee soundness, and what
closed the gap. Section~\ref{sec:discussion} discusses the resulting
governance picture and limitations; Section~\ref{sec:conclusion}
concludes.

\section{Formalization Target and Verified Results}\label{sec:formalization}

% Scope (PLAN.md §7.2): definition of weak solutions, the two capstone theorems,
% exact scope, and the main analytic difficulties.
% Statements MUST match claims/formalization-scope.md (pinned at
% leray-hopf v0.1.0-rc1, commit 7c15710a). No claim beyond that file.

The formalization target is the classical existence theorem for
Leray--Hopf weak solutions of the homogeneous, incompressible
Navier--Stokes equations~\cite{leray1934,hopf1951}: given a
divergence-free initial datum $u_0 \in L^2$, a viscosity $\nu > 0$, and a
time horizon $T > 0$, a weak solution on $[0,T]$ exists. We report a
Lean~4~\cite{moura2021lean4} formalization of this theorem, built on
mathlib~\cite{mathlib2020}, on two domains: the periodic unit 3-torus
$\mathbb{T}^3$ and whole space $\mathbb{R}^3$. All statements below are
quoted from the pinned reference commit of the public repository
\texttt{uda-lab/leray-hopf}, release \texttt{v0.1.0-rc1}, commit
\texttt{7c15710a7b9068a2aa105fc7c11b432e7685b7b5}~\cite{lerayhopf2026repo};
the repository's own scope statement (\texttt{README.md} at that commit)
is treated as the upper bound on what this section may claim.
% CLAIM: CLM-001 (leray-hopf@7c15710a7b9068a2aa105fc7c11b432e7685b7b5)

\subsection{Weak solutions as formalized}

A solution is a proof-carrying structure,
\texttt{LerayHopfSolutionFull} (on $\mathbb{T}^3$) and
\texttt{LerayHopfSolutionFull\_R3} (on $\mathbb{R}^3$), parameterizing a
curve $u : \mathrm{Time} \to D_\sigma$ together with the fields in
Table~\ref{tab:solution-fields}; both are definitional abbreviations of a
single domain-neutral \texttt{Galerkin.LerayHopfSolution} structure shared
by the two lanes. Two restrictions of the formalized notion are worth
flagging explicitly, since they are easy to over-read from the field
names alone. First, \texttt{weak\_eq} is tested only against
\emph{separated-variable} test functions $\psi(t)\!\cdot\! w(x)$, for a
scalar temporal factor $\psi$ and a fixed spatial test vector $w$ drawn
from the domain's divergence-free test class -- not a general space-time
test function $\varphi(t,x)$; the relation between the two formulations
is out of scope, neither proved nor assumed here. Second,
\texttt{initial\_trace} is a one-sided limit $u(t) \to u_0$ in the ambient
$L^2$ norm as $t \to 0^+$, not a weak-continuity statement on all of
$[0,T]$.

\begin{table}[h]
\centering
\caption{Solution-structure fields realizing the weak formulation.}
\label{tab:solution-fields}
\begin{tabular}{@{}p{0.20\linewidth}p{0.70\linewidth}@{}}
\toprule
Field & What is guaranteed \\
\midrule
\texttt{u} & valued in the closed divergence-free subspace $D_\sigma$ at every $t$ (type-level) \\
\texttt{weak\_eq} & weak Navier--Stokes identity against separated-variable tests $\psi(t)w(x)$, no forcing term \\
\texttt{energy\_ineq} & $\tfrac12\|u(t)\|^2 + \int_0^t \mathrm{dissip}(\nu,u(s))\,ds \le \tfrac12\|u_0\|^2$ for $t\in[0,T]$ \\
\texttt{initial\_trace} & $u(t)\to u_0$ in $L^2$ as $t\to 0^+$ (one-sided trace at $t=0$ only) \\
\texttt{energy\_class} & a.e.-in-$t$ membership in the $H^1$ regularity class, plus interval-integrable viscous dissipation \\
\texttt{u\_aestronglyMeasurable} & a.e.-strong measurability of $t \mapsto u(t)$ into the ambient $L^2$ space \\
\bottomrule
\end{tabular}
\end{table}

\subsection{Capstone existence theorems}

The two capstone statements are quoted verbatim from the pinned commit
(proof bodies elided):

\begin{lstlisting}
-- LerayHopf/Torus/GalerkinODECapstone.lean  (T^3)
theorem exists_lerayHopf_torus3 (u₀ : L2Sigma) (ν : ℝ) (hν : 0 < ν)
    (T : ℝ) (hT : 0 < T) :
    ∃ F : Torus3NSForms, Nonempty (LerayHopfSolutionFull F ν T u₀)

-- LerayHopf/R3/GalerkinODECapstone.lean     (R^3)
theorem exists_lerayHopf_r3 (u₀ : L2Sigma_R3) (ν : ℝ) (hν : 0 < ν)
    (T : ℝ) (hT : 0 < T) :
    ∃ (𝔊 : R3GalerkinScheme) (F : R3NSForms 𝔊),
    Nonempty (LerayHopfSolutionFull_R3 𝔊 F ν T u₀)
\end{lstlisting}
% CLAIM: CLM-001 (decl:exists_lerayHopf_torus3, decl:exists_lerayHopf_r3)

Both are unconditional in their stated hypotheses: for \emph{every}
divergence-free $u_0$, viscosity $\nu>0$, and horizon $T>0$, a solution
exists. On $\mathbb{R}^3$ the witness scheme $\mathfrak{G}$ (a concrete
Galerkin scheme built from a Schwartz divergence-free basis) is itself
produced by the proof, not assumed.

\subsection{Scope and non-claims}

The theorems above are exactly what is proved -- no more. This
formalization does \emph{not} claim: smoothness or any regularity beyond
the $H^1$ energy class exhibited by \texttt{energy\_class}; uniqueness or
non-uniqueness of the constructed solution; an equation with an external
forcing term (the homogeneous case only); a single solution valid
uniformly on $[0,\infty)$ (each capstone fixes a finite horizon $T$, taken
as an arbitrary but fixed input); or the weak formulation against general
space-time test functions (only the separated-variable class
$\psi(t)\!\cdot\! w(x)$, as above). A further, purely engineering scope
boundary: the repository contains an explicit
\texttt{LerayHopf.Experimental} namespace holding incomplete
Bochner-space-in-time developments that neither capstone depends on and
that are excluded from the release surface described in
Section~\ref{sec:release-surface} below.

\subsection{Main analytic difficulties}

The 96-file, $\sim$42{,}000-line \texttt{LerayHopf/} tree (measured at the
pinned commit; \texttt{evidence/metrics/formalization-metrics.json})
% CLAIM: CLM-001 (formalization-metrics.json, measured at leray-hopf@7c15710a)
is organized around three domain-neutral layers shared by both lanes, plus
per-domain instantiations. \emph{Galerkin approximation and finite-dimensional ODE
existence}: for each finite-dimensional Galerkin subspace $V_n$, the
truncated quadratic vector field is shown $C^1$ with a dissipative
inner-product estimate, and forward-global existence on $V_n$ is obtained
by a tiling/gluing argument over local Picard--Lindelöf solutions driven
by an a~priori energy bound, rather than by invoking a global
existence theorem from mathlib directly. \emph{Spatial compactness}:
$\mathbb{T}^3$ uses a Fourier mode-tail Rellich argument; $\mathbb{R}^3$
uses a local Rellich--Kondrachov embedding on balls (the one classical
input not otherwise derived from mathlib) combined with a
Fréchet--Kolmogorov criterion and a countable diagonal argument over an
exhausting sequence of balls. \emph{Aubin--Lions compactness in time and
limit passage}: a shared Gelfand-triple/Bochner-space layer supplies the
time-compactness and weak-limit toolkit consumed by both lanes' passage
from the Galerkin sequence to a solution of the weak equation, energy
inequality, and initial trace; the $\mathbb{T}^3$ route additionally
exploits a mode-wise spectral decomposition, while the $\mathbb{R}^3$
route uses Steklov time-averaging. The nonlinear convection term on each
domain is handled by a proof-carrying trilinear extension pinned to the
Schwartz/Galerkin test class, not by a claim of a bounded convection
operator on all of $L^2\times L^2\times L^2$.

\subsection{Kernel-only axiom check}

Both capstones are \emph{kernel-only}: \texttt{\#print axioms} on
\texttt{exists\_lerayHopf\_torus3} and \texttt{exists\_lerayHopf\_r3}
returns exactly the three standard Lean kernel axioms --
\texttt{propext}, \texttt{Classical.choice}, \texttt{Quot.sound} -- with
zero project-declared axioms and no \texttt{sorryAx} on either theorem.
This is pinned by a CI script (\texttt{scripts/check-axioms-live.sh})
that runs the live \texttt{\#print axioms} check and fails on any
unexpected or missing axiom.

\subsection{Release-surface enforcement and build attestation}\label{sec:release-surface}

\texttt{import LerayHopf} is defined as the complete release surface:
both capstones plus every supporting file they depend on. A dedicated
guardrail script, \texttt{scripts/check-release-cone.sh}, walks the
transitive import closure of that entry point and fails if any reachable
file contains a \texttt{sorry}, a project \texttt{axiom}/\texttt{opaque},
or a placeholder namespace; this check runs as a fast, buildless textual
guardrail on every pull request. The \texttt{LerayHopf.Experimental}
namespace mentioned above is, by this same enforcement, unreachable from
the release surface. A separate, heavier step -- a full \texttt{lake
build} together with the live axiom pin -- is run manually, on demand,
against a specific commit, and recorded as a release-candidate build
attestation; it certifies exactly the attested commit SHA, not
continuously the branch tip. The results reported in this section are
attested against commit \texttt{7c15710a} under release tag
\texttt{v0.1.0-rc1} in exactly this manner, not inferred from a
continuously green push-triggered build.

\section{Harness Design for Agent-Assisted Formalization}\label{sec:workflow}

% The proposed harness-architecture figure is tracked in a separate figure
% issue.  Do not add a rendered figure or a figure reference until the
% author-supplied artwork is available.
%
% Historical workflow facts in this section are bound by CLM-007 to
% analysis/workflow-evolution.md and the six coded session indices in
% evidence/session-index/.  The architecture is an end-of-campaign synthesis,
% not a claim that every component was present throughout the campaign.

We reconstruct the campaign's end-state harness from controls introduced at
different stages. It separates scientific authority and implementation from
checks matched to the mathematical statement, Lean artifact, and public
declaration set, together with operational control and evidence capture.
% CLAIM: CLM-009
The resulting architecture is a synthesis of the final project state, not a
claim that every component operated throughout the campaign.

\subsection{Scientific authority and issue-scoped work}

The scientific owner retained authority over the mathematical scope. An issue
specified the intended result and acceptance conditions; for statements that
were easy to overgeneralize, a statement card recorded the hypotheses,
target conclusion, and permitted proof boundary before implementation. The
issue and statement card thus served different purposes: the issue bounded a
unit of work, while the card made the mathematical contract reviewable.
% CLAIM: CLM-009

The orchestrator coordinated work but did not edit Lean files. Implementing
agents received issue-scoped Git worktrees, separate working directories tied
to distinct branches, and returned changes through pull requests. The pull
request consolidated the committed candidate, validation results, and review
record. Coded records show that coordination and implementation were distinct
roles from the earliest observed session; later phases made the pull
request their durable boundary. % CLAIM: CLM-007, EV-2076, EV-0925

For the large refactoring campaign, a read-only scout examined repository
state and issue assumptions before an implementing agent was assigned.
Implementation and inspection therefore did not share edit authority. Two
review roles were also separated: one checked module structure, and another
compared public declarations before and after each refactor. This role
structure was observed in the July campaign; it should not be projected onto
earlier phases. % CLAIM: CLM-007, EV-1669

The minimal change contract was therefore as follows. The issue fixed scope
and acceptance conditions, a worktree held one implementation candidate, and
the pull request recorded the checks applicable to the changed artifact. The
scientific owner revised the statement contract after semantic findings;
implementers addressed code or build failures in the worktree. A candidate
advanced only after the applicable check results and a reviewer-attributed
completion record were attached to the pull request. % CLAIM: CLM-009

\subsection{Verification by artifact}

The harness assigned a distinct check to each artifact at risk.
% CLAIM: CLM-009

\paragraph{Mathematical statement.}
An agent independent of the implementer examined whether the proposition
matched the intended mathematics and whether its hypotheses were sufficient.
Reviewers tested candidate statements with concrete counterexamples when
appropriate. This semantic gate assessed the intended proposition and returned
scope corrections to the statement contract; the contrasting outcomes are
analyzed in Section~\ref{sec:incidents}.
% CLAIM: CLM-005, CLM-007, EV-0925

\paragraph{Lean artifact and axiom footprint.}
The build checked elaboration and type correctness. A separate
\texttt{\#print axioms} result recorded the assumptions of designated
declarations, while release-surface checks excluded experimental modules.
Together they certified the Lean artifact and its dependency boundary. A
release candidate then received a recorded full-build attestation for its
specific commit.
% CLAIM: CLM-009

\paragraph{Public declarations.}
Refactoring was checked against the prior public declaration set, independently
of downstream use. The declaration-level comparison certified interface
preservation by comparing the candidate inventory and, where exact
preservation was required, statement text. Section~\ref{sec:incidents}
describes the omission that motivated this boundary.
% CLAIM: CLM-007, EV-1669

\paragraph{Independent review record.}
A pull request required durable evidence that a reviewer distinct from the
implementer had examined the change. A dispatched review did not satisfy the
gate. The pull request needed an attributable record naming the reviewer and
result; the gate checked this completion evidence separately from the
substantive findings.
% CLAIM: CLM-007, EV-1669

\subsection{Ownership, liveness, and resource control}

Parallel work used independently named worktrees with one active implementing
owner per worktree. Shared import files were assigned to a single writer.
Silence alone did not transfer ownership: the orchestrator was expected to
check agent and process state before replacing an unresponsive worker. These
rules address two different risks, concurrent edits to shared state and
duplicated work caused by an incorrect liveness judgment.
% CLAIM: CLM-007, EV-2076, EV-1669, EV-0247

Lean builds shared a cache and ran under a global lock. Before a build, the
harness checked available memory across the container and removed stale helper
processes; under pressure it reduced Lean's thread count. These checks formed
the final operating arrangement and determined whether verification could
complete. Assurance about the mathematical statement came from the
checks above. % CLAIM: CLM-007, CLM-009, EV-0247

Agent assignments changed across the campaign. The durable role definitions
were therefore expressed through permissions and recorded outputs rather than
through model names. % CLAIM: CLM-007, EV-2076, EV-0925, EV-1669

\subsection{Evidence capture and observed events}

Session records, pull-request discussions, build results, and release
artifacts were retained as separate evidence sources. Claim identifiers link
paper statements to those sources, and reconstructed intervals are marked
separately from primary session evidence. This evidence layer supports audit
and later reconstruction alongside the verification artifacts.
% CLAIM: CLM-009

Across the campaign period from 2026-06-10 through the 2026-07-20 release,
the collected session logs record at least 15.1M output tokens, an
API-equivalent cost of \$6{,}973, and 199.1h of active engagement under a
5-minute inter-event gap cap. API-equivalent cost applies the published API
rates to recorded token classes; active engagement sums adjacent events only
when their separation is at most five minutes. The values are lower bounds
because the evidence store has known coverage gaps. % CLAIM: CLM-002

For the Vertex-hosted completion phase on July 2--3, 2026 (Pacific Time),
billing reconciliation kept three quantities separate. The provider export
measured a machine-wide actual charge of JPY~54{,}868; applying published
prices to the project logs yielded an unconditional calculation of
\(\$258.28\). Under the stated interpretation of the provider's aggregate
token metric, the project's derived share of the actual charge was
JPY~41{,}421--41{,}796. This attribution range is conditional on that metric
interpretation. % CLAIM: CLM-003

Six sessions were selected purposively for reconstruction value and coded by
the incident analyst against a fixed event vocabulary, separately from
evaluation of the formal result. Table~\ref{tab:event-codes}
describes only those sessions. The event families connect the evidence to the
design: statement and mathematical events call for semantic gates, whereas
ownership, handoff, and resource events call for operational and
review-completion controls.
% CLAIM: CLM-004

\begin{table}[t]
\centering
\caption{Observed events in the six purposively selected sessions.}\label{tab:event-codes}
\begin{tabular}{@{}p{0.72\linewidth}r@{}}
\toprule
Observed event & Count \\
\midrule
Statement mismatch or weakening & 4 \\
Recovery action & 19 \\
Resource failure & 12 \\
Worktree ownership conflict & 5 \\
Uncorrected false success report & 0 \\
\bottomrule
\end{tabular}
\end{table}

No uncorrected false success report, defined here as a mismatch
between a runtime success report and on-disk state, was observed in the coded
sessions. This is a descriptive count for the purposive sample.
% CLAIM: CLM-004

\section{Incidents and Recovery}\label{sec:incidents}

A successful Lean build establishes that the declarations in the checked
dependency graph are well typed. The four cases below concern boundaries it
does not reach: whether a theorem states the intended mathematics, whether
a refactor preserved every public declaration, and whether the runtime can
complete the next build. The evidence combines repository records with six
agent sessions selected for their reconstruction value, neither randomly
nor exhaustively, so the cases document observed failures and
recoveries rather than failure rates. PR numbers refer to the leray-hopf
repository~\cite{lerayhopf2026repo}; the cases do not depend on them.
% CLAIM: CLM-004, CLM-008

\subsection{Statement correctness}\label{sec:inc001}

\emph{Case A. A false generalization detected before release.}
The limit passage takes place in an evolution triple
$V \hookrightarrow H \hookrightarrow V'$; in the classical theory, $H$ is
the divergence-free $L^2$ space and $V$ its $H^1$ subspace. A standard
lemma states that if $u \in L^2(0,T;V)$ and its weak time derivative
satisfies $u' \in L^2(0,T;V')$, then $u$ agrees almost everywhere with a
function continuous into $H$; classically this is what gives a weak
solution a value at every time. The exponents matter because the proof
integrates the pairing $t \mapsto \langle u'(t), u(t)\rangle$ over time,
which Cauchy--Schwarz makes integrable when both exponents equal~2. A
public lemma in the project, declared over an abstract triple, asserted
the conclusion for arbitrary exponents $p,q\geq 1$, with
$u \in L^p(0,T;V)$ and $u' \in L^q(0,T;V')$, although the only proof
route examined was the case $p=q=2$. The declaration
carried an explicit \texttt{sorry}, and a later proof attempt checked only
that no capstone theorem depended on it. The false general statement
remained public for about a month and was repeated in project
documentation. % CLAIM: CLM-008, INC-001

\emph{Detection and response.}
During release preparation, a statement reviewer tested the boundary case
$p=q=1$. Take $H=\ell^2$ with a weighted triple in which the $n$th basis
vector has $V$-norm $2^n$ and $V'$-norm $2^{-n}$, and let $u$ consist of
spikes of unit $H$-norm along successive basis vectors, placed on
disjoint intervals of width $4^{-n}$ that accumulate at $t=0$. The widths
$4^{-n}$ against the $V$-norms $2^n$ make $u \in L^1(0,T;V)$, and the
$V'$-norms $2^{-n}$ make $u' \in L^1(0,T;V')$, so the stated hypotheses
hold, but the unit peaks accumulating at a point admit no
almost everywhere equal representative continuous into $H$. The problem
was therefore the proposition, not a difficult proof. The repair restricted
the declaration to $p=q=2$ and moved it outside the public imports
(PR~\#162). A follow-up added a statement card recording the exact type,
the literature reference, and the boundary cases checked, together with a
regression check that fails if the unrestricted signature reappears or a
placeholder lacks a card (PR~\#170).
% CLAIM: CLM-008, INC-001

\emph{Lesson.}
A placeholder records only that a proof is absent; it provides no evidence
that the proposition is true. Before proof search, reviewers should test the
boundary cases admitted by a public statement, especially when its
quantifiers or exponents exceed those used in the motivating argument.
% CLAIM: CLM-008, INC-001

\emph{Case B. A missing hypothesis found before merge.}
On $\mathbb{T}^3$ the convection term enters the Galerkin equations
through the trilinear form $b(u,v,w) = \int (u\cdot\nabla v)\cdot w$, and
the energy inequality rests on its antisymmetry $b(u,v,w) = -b(u,w,v)$ for
divergence-free $u$. The Galerkin scheme truncates Fourier modes to a
finite box, and the truncated form $b_n$ is a finite sum over pairs of
modes in the box. A second placeholder asserted antisymmetry of $b_n$ for
arbitrary $v$ and $w$. The standard proof, however, reindexes the
inner sum by an involution of the modes, which stays inside the box only
when the Fourier coefficients of $v$ and $w$ vanish outside it, that is,
when both lie in the finite projection space $V_n$. Two proof revisions had
attempted the unrestricted statement before an independent reviewer
identified this missing closure condition. % CLAIM: CLM-008, INC-002

\emph{Detection and response.}
The reviewer tested inputs outside $V_n$. A numerical diagnostic gave a
discrepancy between $b_n(u,v,w)$ and $-b_n(u,w,v)$ of approximately $80.7$
without the projection hypotheses and approximately $10^{-13}$ after
imposing them. These values corroborated the structural diagnosis; they
were not a certified counterexample. The repair added the two hypotheses
$P_n v = v$ and $P_n w = w$, where $P_n$ is the projection onto $V_n$,
completed the Lean proof, and removed the placeholder before merge
(PR~\#27). % CLAIM: CLM-008, INC-002

\emph{Lesson.}
For identities derived from truncation or projection, review should test
inputs outside the invariant subspace. Case B shows a defect caught by that
review before merge, whereas the analogous review in Case A came only at
release preparation. % CLAIM: CLM-005, CLM-008, INC-001, INC-002

\subsection{Public declarations during refactoring}

The construction of the convection form on $\mathbb{T}^3$ and on
$\mathbb{R}^3$ shared a layer of tensor and gluing lemmas that had been
written twice, once per domain. A refactor replaced both copies with
instances of a single parametrized module and was required to preserve
every public theorem of each domain file, most as thin instantiations of
the shared one (PR~\#120). The refactored code built successfully, but one
public theorem on the $\mathbb{R}^3$ side, which had no downstream users,
was omitted from the list of theorems to reconnect and disappeared, so
neither the Lean build nor an earlier review of module structure revealed
the omission. % CLAIM: CLM-008, INC-004

\emph{Detection and response.}
A separate reviewer compared the public declaration inventory before and
after the refactor, statement by statement. The missing theorem was
restored as a thin wrapper around the shared implementation with its
original statement, and the reviewer repeated the comparison before merge.
% CLAIM: CLM-008, INC-004

\emph{Lesson.}
When a refactor promises interface preservation, the review target must be
the complete public declaration set, including declarations with no
current users. A build tests dependencies. An inventory comparison tests
preservation. % CLAIM: CLM-008, INC-004

\subsection{Resource availability and orchestration}\label{sec:inc005}

The VPS container had an effective memory limit of 3.42\,GiB and no swap.
Review jobs from the previous night left fourteen helper process chains
running after their Git worktrees were inactive, consuming about
1--1.5\,GiB. At the same time, the orchestrator misclassified a slow agent
as inactive and assigned a replacement to the same Git worktree. The
replacement did not receive the shared cache rule and started a full
mathlib build. Under this load, Lean builds repeatedly ended with
exit code 137 and no compiler diagnostic. % CLAIM: CLM-008, INC-005

\emph{Detection and response.}
Monitoring had covered the visible Lean process rather than memory
available to the entire container. Roughly seventy minutes after an earlier
status check had reported the container as healthy on that basis, the
author's next status request reached the session only after an eight-minute
delay, and was followed at once by the author's explicit instruction to
kill stale processes. Only then did a memory check report 606\,MiB
available, and inspection of each process working directory identified the
fourteen stale chains. A cumulative kill counter maintained by the Linux
control group (cgroup) subsystem, read during the response, was taken at
the time as confirmation that the out of memory (OOM) killer had terminated
the failed builds; an audit of host kernel logs later overturned that
reading, because every counted kill belonged to a host-level OOM event two
days earlier and no kernel OOM occurred on the incident day. The mechanism
that terminated the builds remains unidentified. Removing the stale
processes raised available memory to 2.2\,GiB. The worktree returned to the
original agent, the duplicate stopped, and the build completed under a
global lock. The three PRs blocked by the failed builds merged later that
day (PRs~\#174--\#176).
% CLAIM: CLM-008, INC-005

\emph{Lesson.}
An orchestrator must observe the complete process tree and available
memory, not only the foreground prover. After the incident, helper-process
cleanup, cache reuse, build serialization, and a pre-build memory check
entered the standing instructions and dispatch prompts. Once
monitoring tracked memory
available to the whole container, it also detected a later peak caused by
one legitimate build, so measuring the resource directly is more general
than cleanup tailored to one cause. The
corrected diagnosis adds an evidence rule: a causal account assembled only
from observations inside the agent session can misattribute the mechanism,
so it must be checked against records outside the session.
% CLAIM: CLM-008, INC-005

\section{Engineering Implications and Limitations}\label{sec:discussion}

\subsection{Recommendations}

The principal design question is which artifact each control can assess.
The observed cases support the following recommendations for projects in
which agents develop a substantial formalization under a mathematical
author. % CLAIM: CLM-005, CLM-009

\emph{Keep mathematical authority outside the implementation role.}
Elaboration establishes that a proposition is well formed, not that it
expresses the intended mathematics. The author should approve a statement
card before implementation when the scope is delicate, and a reviewer who
did not write the proof should test its hypotheses and boundary cases, as
Cases A and B illustrate. % CLAIM: CLM-005, INC-001, INC-002

\emph{Use the GitHub workflow to route mathematical feedback.}
The mathematical value of \texttt{github-driven-workflow} came from its
routing rule: semantic findings return to the author and the issue
contract, implementation findings return to the worktree, and the PR
preserves the candidate, the applicable checks, and the decision in one
durable record. % CLAIM: CLM-007, CLM-009, EV-2076, EV-1669

\emph{Separate observation from modification.}
A read-only scout inspected repository assumptions, while specialist
reviewers examined mathematical scope, module structure, and the public
declaration inventory. The declaration comparison in PR~\#120 found a
missing theorem after both the build and a structural review had passed.
% CLAIM: CLM-007, CLM-008, CLM-009, INC-004

\emph{Place resource rules in the dispatch path.}
The VPS made resource contention visible, but the same mechanism applies to
laptops and workstations that run many agents continuously: abandoned agent
and helper process trees retain memory, while concurrent Lean builds add
large transient demand. Dispatch rules should therefore clean up those
process trees, bound agent and build concurrency, reuse the build cache,
serialize heavy builds, and check the memory available to the whole machine
or container rather than to the foreground prover alone. Recording such a
rule in an agent's memory is insufficient because later agents and tools
change; standing instructions and dispatch prompts carry it to replacement
agents and later sessions.
% CLAIM: CLM-007, CLM-008, CLM-009, INC-005

\subsection{Evidence and limits}

This observational account evaluates one author-supervised formalization
workflow as a whole, including the author's own interventions. Comparative
evaluation across projects would be needed to isolate the contribution of
an individual control and to assess how far it transfers.
% CLAIM: CLM-004, CLM-005

Early session records are incomplete, and the models and proprietary tools
changed during the project outside its control. Reconstructed events are
distinguished from primary session evidence and carry an explicit
confidence assessment. The recommendations are therefore expressed through
durable artifacts, assigned responsibilities, and observable checks rather
than a requirement to use the same software. % CLAIM: CLM-008, INC-003

The paper repository provides redacted evidence excerpts,
claim-to-evidence manifests, and scripts for the reported metrics. Raw
session logs remain private because they contain information about the
people and systems recorded in them.

\section{Conclusion}\label{sec:conclusion}

We formalized the classical existence theorem for Leray--Hopf weak
solutions on $\mathbb{T}^3$ and $\mathbb{R}^3$ in Lean~4 and mathlib. At
the pinned release, the two capstone theorems depend on no project axioms.
% CLAIM: CLM-001

The engineering contribution is an observed division of responsibility.
The author controls theorem scope, while an orchestrator assigns bounded
work through task contracts, isolated worktrees, and reviewed PRs.
Specialist reviewers assess the statement, Lean assumptions, and public
declarations. Standing dispatch rules govern the resources shared by
agents and Lean builds. The incidents show defects observed at each of
these boundaries in one supervised project, and the retained evidence
records what was checked. % CLAIM: CLM-007, CLM-009


\bibliographystyle{IEEEtran}
\bibliography{references}

\end{document}
